\section{Discussion}
\label{sec:discussion}

\subsection{Role of the Latent Representation and Group-Aware Augmentation}

The results support the hypothesis that self-supervised representation learning can guide unsupervised feature selection for IIoT traffic. \rev{The controlled comparisons isolate the contribution of the evaluation space: at identical feature budgets, subsets chosen through the frozen VICReg encoder yield higher downstream macro-F1 than subsets chosen by the same search in the standardized input space, and the latent search also produces markedly more stable subsets across folds (Jaccard similarity 0.69 versus 0.40). The latent space therefore provides a more informative and more reproducible criterion for joint subset evaluation than raw-space clustering.}

VICReg is well suited to this role because it does not require negative pairs and its variance term penalizes collapsed representations. \rev{Within this objective, group-aware masking and conventional independent masking at the same rate produced comparable downstream performance. Group-aware masking retains its design advantage of keeping each semantic measurement group internally consistent in the augmented views, and it provides a natural mechanism to encode domain knowledge about which measurements belong together; the Context Groups were rebuilt for two additional schemas without modifying the method.} ContrastiveXAI-FS continues to select original features, so the operational meaning of the selected inputs is preserved while interactions among timing, packet-rate, addressing, header, fragmentation, and log-derived measurements influence the subset objective.

\subsection{Operating Region and Feature Budgets}

\rev{At its natural operating point, the method retains on average 59 of 71 features (61 in the descriptive configuration) and remains within $\pm$0.01 macro-F1 of the complete input across three master seeds. Under very restrictive budgets, MCFS is the most stable method (0.881 at $k = 10$), whereas \method{} reaches 0.536 at $k = 10$ and 0.878 at $k = 25$. This establishes the operating region of the method:} the latent objective is designed to remove attributes that do not improve the learned behavioral organization, not to enforce a small cardinality. When a budget of ten features is required, the encoder, trained on complete inputs with at most two masked groups, evaluates inputs in which most dimensions are masked, and an internal Silhouette criterion is a weaker surrogate for class separation in that regime. Exposing the encoder to masking levels closer to the intended budget, or combining latent preselection with a budget-aware refinement stage, are natural extensions for compact deployments.

\subsection{Implications for IIoT Security Practice}

\rev{The practical value of retaining most of the 71 features lies in three properties rather than in computational savings, which were not measured.} First, the refinement removes attributes that do not support the learned behavioral structure, chiefly the inter-packet timing statistics, which are removed in most folds and are known to reflect device polling and reporting cycles; this yields a cleaner input without loss of detection performance relative to the complete feature set. Second, the selected subset is the input of the cluster explanations, so the analyst receives behavioral profiles expressed in a validated set of original measurements. Third, the procedure is entirely label-free and can therefore be applied before comprehensive attack labels exist. When the main requirement is the smallest possible number of monitored attributes, compact classical selectors such as MCFS are preferable; when the priority is to preserve unlabeled traffic structure and support analyst-facing explanations, \method{} provides a suitable operating point.

The SHAP profiles translate latent clusters into original traffic features, such as TCP termination flags, TTL levels, TTL variability, and payload sizes. These profiles are investigative cues rather than attack diagnoses: the same measurements can also arise from legitimate device, routing, or topology differences, and they are intended to support analyst prioritization and follow-up investigation.

\subsection{Computational Cost and Scalability}
\label{sec:cost}

ContrastiveXAI-FS is designed as an offline feature-selection pipeline. Its principal computational stages are VICReg training and the bidirectional subset search, which are repeated independently for each cross-validation fold in the present evaluation. Every previously unseen candidate subset requires encoder inference, $k$-Means clustering, and a Silhouette estimate, so the search is more demanding than a filter that assigns one score to each feature; the subset cache prevents repeated evaluation within a fold. \rev{Because the objective is estimated on a fixed training-fold subsample, the cost per candidate is independent of the training-fold size, and the method was executed on a CPU-only notebook for all three datasets, including the 115-feature N-BaIoT schema.} Once a subset and downstream classifier have been fixed, an operational IDS collects the retained attributes and applies only the deployed classifier; runtime, memory, throughput, and energy of deployment were not measured, and no efficiency claim is made.

\subsection{External Validity and Deployment Considerations}

\rev{The Context Group principle transferred to two schemas that differ substantially from DataSense: per-packet damped-window statistics in N-BaIoT and bidirectional flow records in WUSTL-IIoT-2021. On N-BaIoT, the method was statistically equivalent to the complete feature set and superior to MCFS, reproducing the DataSense pattern; on WUSTL-IIoT-2021, it was competitive in binary detection with a more variable subset size. The number and composition of the groups remain dataset dependent, and the selected features, cluster structure, and SHAP profiles are specific to each environment.}

\rev{The main experiments use sample-level folds and therefore measure within-dataset IID generalization. The grouped protocols quantify the effect of keeping related traffic out of the test partition: when complete attack executions and benign time blocks are held out, binary detection is unchanged, and when complete devices are held out, binary detection remains above 0.90 while the eight-class task becomes harder because several attack categories are recorded on few devices. In both grouped protocols, none of the differences between \method{} and the complete feature set is statistically significant.}

The current framework is offline. Under concept drift, changes in devices, protocols, workloads, or attacker behavior may reduce the relevance of the learned representation and explanations, so a deployment would monitor feature distributions, detection performance, or cluster stability and repeat the offline procedure when meaningful changes are detected. \rev{The cross-pipeline cluster-stability measure introduced in this revision provides a direct quantity for such monitoring.}

\subsection{Scope of the Explainability Results}

\rev{The explanations are computed for a Random Forest proxy whose agreement with the latent partition was measured on data unseen by every stage of the pipeline (0.97 balanced accuracy on the outer test folds), and the partition itself is reproducible across folds and master seeds (ARI 0.94). The clusters chosen for detailed interpretation are selected by a quantitative clarity index computed for all clusters, which combines held-out proxy fidelity and profile concentration. The candidate cluster range bounds the granularity of the explanation; within that range, the latent Silhouette operates on the plateau of its profile.} Context-Group comparisons rely on feature-level contributions because raw group sums would favor larger groups.
